% GENERATED FILE. Edit canonical lessons, then run npm run generate:books.
% canonical-source-hash: fnv1a64:6f4fab98e1d7a7f0
% canonical-lessons: RU-C493-vlyublyonnyy, RU-C493-potseluy, RU-C493-flirtovat, RU-C493-nezhnost, RU-C493-laskovyy

\chapter{Sweethearts}
\label{ch:ru-sweethearts}
\hlchaptermodality{493}

\begin{chapteropening}
\textbf{By the end of this chapter:} \emph{I can talk about being in love and showing affection.}

Complete the last lesson of chapter 493: I can talk about being in love and showing affection.
\end{chapteropening}

\section[\texorpdfstring{\ru{влюблённый} (vlyublyónnyy)}{влюблённый (vlyublyónnyy)}]{\ru{влюблённый} (vlyublyónnyy) — in love; a lover}

\label{lesson:RU-C493-vlyublyonnyy}



\begin{quote}
Before the new one: say the Russian for twins, then the Russian for related as cousins.
\end{quote}

\subsection*{You'll want to know: \ru{влюблённый}}
\textbf{\ru{влюблённый}} (\emph{vlyublyónnyy}) — \textquotedblleft{}in love; a lover\textquotedblright{}.

\begin{quote}
\ru{Влюблённые} \ru{гуляли} \ru{по} \ru{парку} \ru{до} \ru{утра}.
\end{quote}

\begin{quote}
\emph{Vlyublyónnyye gulyáli po párku do utrá.}
\end{quote}

\begin{quote}
The lovers walked in the park until morning.
\end{quote}

From \emph{\ru{влюбиться}} — to fall in love. \emph{\ru{Влюблён} \ru{в} \ru{кого}} — in love with someone.

\begin{cousinweb}
\emph{\ru{влюблённый}} is the past participle of \emph{\ru{влюбиться}} 'to fall in love', built on \emph{\ru{в}} 'into' and \emph{\ru{любить}} 'to love'; EtymDB lists the participle without a source.
\end{cousinweb}

\begin{grammarlens}[title={What you've built}]
One more word for telling, explaining and describing at B1.
\end{grammarlens}

\subsection*{Your turn}
\begin{itemize}
\raggedright
  \item \emph{Say it:} \emph{vlyublyónnyy}
  \item \emph{Say it:} \emph{vlyublyónnyy}, once more
  \item \emph{Say it:} a sentence of your own with \emph{vlyublyónnyy}
  \item \emph{Recall:} say the Russian for twins, then the Russian for related as cousins, then say \emph{vlyublyónnyy} again
\end{itemize}

\begin{tcolorbox}[breakable,colback=teal!4,colframe=teal!35!black,title={Before you move on}]
What does \ru{влюблённый} mean? (\textquotedblleft{}In love; a lover\textquotedblright{}.) Say it once more.
\end{tcolorbox}
\section[\texorpdfstring{\ru{поцелуй} (potselúy)}{поцелуй (potselúy)}]{\ru{поцелуй} (potselúy) — a kiss}

\label{lesson:RU-C493-potseluy}



\begin{quote}
Before the new one: say the Russian for related as cousins, then the Russian for in love.
\end{quote}

\subsection*{You'll want to know: \ru{поцелуй}}
\textbf{\ru{поцелуй}} (\emph{potselúy}) — \textquotedblleft{}a kiss\textquotedblright{}.

\begin{quote}
\ru{Она} \ru{на} \ru{всю} \ru{жизнь} \ru{запомнила} \ru{первый} \ru{поцелуй}.
\end{quote}

\begin{quote}
\emph{Oná na vsyu zhizn zapómnila pérvyy potselúy.}
\end{quote}

\begin{quote}
She remembered her first kiss for the rest of her life.
\end{quote}

Masculine; genitive \ru{поцелуя}. The verb is \emph{\ru{целовать}} — to kiss.

\begin{cousinweb}
\emph{\ru{поцелуй}} is built on \emph{\ru{целовать}} 'to kiss', which EtymDB traces to Proto-Slavic \emph{cělovati} 'to greet; to kiss', derived from \emph{cěl\ru{ъ}} 'whole', as if a kiss once wished someone good health; EtymDB has no entry for \emph{\ru{поцелуй}}.
\end{cousinweb}

\begin{grammarlens}[title={What you've built}]
One more word for telling, explaining and describing at B1.
\end{grammarlens}

\subsection*{Your turn}
\begin{itemize}
\raggedright
  \item \emph{Say it:} \emph{potselúy}
  \item \emph{Say it:} \emph{potselúy}, once more
  \item \emph{Say it:} a sentence of your own with \emph{potselúy}
  \item \emph{Recall:} say the Russian for related as cousins, then the Russian for in love, then say \emph{potselúy} again
\end{itemize}

\begin{tcolorbox}[breakable,colback=teal!4,colframe=teal!35!black,title={Before you move on}]
What does \ru{поцелуй} mean? (\textquotedblleft{}A kiss\textquotedblright{}.) Say it once more.
\end{tcolorbox}
\section[\texorpdfstring{\ru{флиртовать} (flirtovát)}{флиртовать (flirtovát)}]{\ru{флиртовать} (flirtovát) — to flirt}

\label{lesson:RU-C493-flirtovat}



\begin{quote}
Before the new one: say the Russian for in love, then the Russian for a kiss.
\end{quote}

\subsection*{You'll want to know: \ru{флиртовать}}
\textbf{\ru{флиртовать}} (\emph{flirtovát}) — \textquotedblleft{}to flirt\textquotedblright{}.

\begin{quote}
\ru{На} \ru{вечеринке} \ru{он} \ru{флиртовал} \ru{со} \ru{всеми}.
\end{quote}

\begin{quote}
\emph{Na vecherínke on flirtovál so vsémi.}
\end{quote}

\begin{quote}
At the party he flirted with everyone.
\end{quote}

Imperfective; present \ru{флиртую}. \emph{\ru{Флиртовать} \ru{с} \ru{кем}}.

\begin{cousinweb}
\emph{\ru{флиртовать}} is built on \emph{\ru{флирт}}, which was borrowed from English \emph{flirt} (EtymDB), with the verb ending \emph{-\ru{овать}}; EtymDB has no entry for the verb.
\end{cousinweb}

\begin{grammarlens}[title={What you've built}]
One more word for telling, explaining and describing at B1.
\end{grammarlens}

\subsection*{Your turn}
\begin{itemize}
\raggedright
  \item \emph{Say it:} \emph{flirtovát}
  \item \emph{Say it:} \emph{flirtovát}, once more
  \item \emph{Say it:} a sentence of your own with \emph{flirtovát}
  \item \emph{Recall:} say the Russian for in love, then the Russian for a kiss, then say \emph{flirtovát} again
\end{itemize}

\begin{tcolorbox}[breakable,colback=teal!4,colframe=teal!35!black,title={Before you move on}]
What does \ru{флиртовать} mean? (\textquotedblleft{}To flirt\textquotedblright{}.) Say it once more.
\end{tcolorbox}
\section[\texorpdfstring{\ru{нежность} (nézhnost)}{нежность (nézhnost)}]{\ru{нежность} (nézhnost) — tenderness}

\label{lesson:RU-C493-nezhnost}



\begin{quote}
Before the new one: say the Russian for a kiss, then the Russian for to flirt.
\end{quote}

\subsection*{You'll want to know: \ru{нежность}}
\textbf{\ru{нежность}} (\emph{nézhnost}) — \textquotedblleft{}tenderness\textquotedblright{}.

\begin{quote}
\ru{Бабушка} \ru{смотрела} \ru{на} \ru{внука} \ru{с} \ru{нежностью}.
\end{quote}

\begin{quote}
\emph{Bábushka smotréla na vnúka s nézhnostyu.}
\end{quote}

\begin{quote}
Grandma gazed at her grandson with tenderness.
\end{quote}

Feminine. Built on \emph{\ru{нежный}} — tender, gentle.

\begin{cousinweb}
\emph{\ru{нежность}} is built on \emph{\ru{нежный}} 'tender' with \emph{-\ru{ость}}; EtymDB has no entry, and \emph{\ru{нежный}} is usually traced to Proto-Slavic \emph{něž\ru{ь}n\ru{ъ}}.
\end{cousinweb}

\begin{grammarlens}[title={What you've built}]
One more word for telling, explaining and describing at B1.
\end{grammarlens}

\subsection*{Your turn}
\begin{itemize}
\raggedright
  \item \emph{Say it:} \emph{nézhnost}
  \item \emph{Say it:} \emph{nézhnost}, once more
  \item \emph{Say it:} a sentence of your own with \emph{nézhnost}
  \item \emph{Recall:} say the Russian for a kiss, then the Russian for to flirt, then say \emph{nézhnost} again
\end{itemize}

\begin{tcolorbox}[breakable,colback=teal!4,colframe=teal!35!black,title={Before you move on}]
What does \ru{нежность} mean? (\textquotedblleft{}Tenderness\textquotedblright{}.) Say it once more.
\end{tcolorbox}
\section[\texorpdfstring{\ru{ласковый} (láskovyy)}{ласковый (láskovyy)}]{\ru{ласковый} (láskovyy) — affectionate}

\label{lesson:RU-C493-laskovyy}



\begin{quote}
Before the new one: say the Russian for to flirt, then the Russian for tenderness.
\end{quote}

\subsection*{You'll want to know: \ru{ласковый}}
\textbf{\ru{ласковый}} (\emph{láskovyy}) — \textquotedblleft{}affectionate\textquotedblright{}.

\begin{quote}
\ru{У} \ru{нашей} \ru{кошки} \ru{ласковый} \ru{характер}.
\end{quote}

\begin{quote}
\emph{U náshey kóshki láskovyy kharákter.}
\end{quote}

\begin{quote}
Our cat has an affectionate nature.
\end{quote}

Built on \emph{\ru{ласка}} — a caress.

\begin{cousinweb}
\emph{\ru{ласковый}} is built on \emph{\ru{ласка}} 'a caress, kindness'; EtymDB lists \emph{\ru{ласковый}} without a source.
\end{cousinweb}

\begin{grammarlens}[title={What you've built}]
One more word for telling, explaining and describing at B1.
\end{grammarlens}

\subsection*{Your turn}
\begin{itemize}
\raggedright
  \item \emph{Say it:} \emph{láskovyy}
  \item \emph{Say it:} \emph{láskovyy}, once more
  \item \emph{Say it:} a sentence of your own with \emph{láskovyy}
  \item \emph{Recall:} say the Russian for to flirt, then the Russian for tenderness, then say \emph{láskovyy} again
\end{itemize}

\begin{tcolorbox}[breakable,colback=teal!4,colframe=teal!35!black,title={Before you move on}]
What does \ru{ласковый} mean? (\textquotedblleft{}Affectionate\textquotedblright{}.) Say it once more.
\end{tcolorbox}
